\chapter{AES, l’algoritmo di cifratura usato da TCFS}
Nel corso dell'implementazione di TCFS sono stati implementati diversi algoritmi di AES per la cifratura dei file, sono stati valutati CBC, GCM, CTR, e parallelamente, l’AES 256 CBC per la cifratura dei nomi delle cartelle e i nomi dei file. Questo capitolo è dedicato all’analisi delle decisioni progettuali che hanno guidato la selezione di questi algoritmi crittografici e alla discussione dettagliata della loro implementazione all’interno di TCFS. L’obiettivo è di fornire una comprensione approfondita delle motivazioni dietro queste scelte e di come esse contribuiscono a migliorare la sicurezza e l’efficienza del nostro sistema.

\section{Definizione e caratteristiche di AES}
L' \textbf{Advanced Encryption Standard (AES)} è un algoritmo di cifratura a blocchi utilizzato per crittografare dati sensibili. Le sue principali caratteristiche includono:
\begin{itemize}
    \item \textbf{Robustezza}: AES è noto per la sua resistenza agli attacchi crittografici.
    \item \textbf{Efficienza}: L'algoritmo AES è progettato per essere veloce ed efficiente.
    \item \textbf{Flessibilità}: Supporta chiavi di diversi bit (128, 192, 256 bit) offrendo flessibilità nella sicurezza.
    \item \textbf{Standardizzazione}: Essendo uno standard internazionale, AES è ampiamente supportato.
\end{itemize}

\subsection{AES CBC GCM e CTR}
In questa sezione analizzeremo le differenze tra AES CBC, GCM e CTR, i loro vantaggi e svantaggi, e le loro applicazioni per giustificare la scelta implementativa di diffirenziare l'algoritmo di crifratura per i nomi delle cartelle e dei file.

\subsubsection{AES CBC}
AES CBC è un modo di operazione che usa una catena di blocchi per cifrare i dati. Ogni blocco di testo in chiaro viene sottoposto a un’operazione di XOR (OR esclusivo) con il blocco di testo cifrato precedente, o con un vettore di inizializzazione (IV) nel caso del primo blocco, prima di essere cifrato con la chiave AES. Questo modo garantisce che due blocchi uguali di testo in chiaro non producano lo stesso blocco di testo cifrato, aumentando la sicurezza. Tuttavia, AES CBC ha anche alcuni svantaggi:

\begin{itemize} 
\item AES CBC non è parallelizzabile, cioè non è possibile cifrare o decifrare più blocchi contemporaneamente, perché ogni blocco dipende dal precedente. Questo riduce le prestazioni, soprattutto con dati di grandi dimensioni. 
\item AES CBC non fornisce autenticazione, cioè non è possibile verificare se il testo cifrato è stato alterato o manipolato da un attaccante. Si potrebbe usare AES CBC in combinazione con un MAC (Message Authentication Code), che è un codice che garantisce l’integrità dei dati. Questo richiede però una chiave aggiuntiva e un’operazione in più, aumentando la complessità e il costo computazionale. \end{itemize}
AES CBC è adatto per cifrare file locali che non richiedono accesso casuale, cioè che vengono letti o scritti sequenzialmente.

\begin{figure}[h]
    \centering
    \includesvg[width=1\textwidth]{capitoli/risorse/AES_CBC_diagram.svg}
    \caption{Architettura AES CBC}
    \label{fig:AES-CBC-diagramma}
\end{figure}

\subsubsection{AES GCM}
AES GCM è un modo di operazione che usa una modalità contatore per cifrare i dati e una funzione di autenticazione basata sul campo di Galois per garantire l’integrità dei dati. Ogni blocco di testo in chiaro viene sottoposto a un’operazione di XOR con il risultato della cifratura di un contatore, che è un numero che incrementa ad ogni blocco. Il contatore viene inizializzato con un nonce (numero usato una sola volta) e un IV. Il testo cifrato viene poi usato per calcolare un tag di autenticazione, che è un codice che permette di verificare se il testo cifrato è stato alterato o manipolato da un attaccante, il tag in GCM è nella pratica il valore del \textbf{polinomio di Galois}. AES GCM ha diversi vantaggi rispetto a AES CBC:

\begin{itemize} 
\item AES GCM è parallelizzabile, cioè è possibile cifrare o decifrare più blocchi contemporaneamente, perché ogni blocco dipende solo dal contatore. Questo aumenta le prestazioni, soprattutto con dati di grandi dimensioni. 
\item AES GCM fornisce autenticazione, cioè è possibile verificare se il testo cifrato è stato alterato o manipolato da un attaccante. Questo aumenta la sicurezza e la robustezza dell’algoritmo. Inoltre, AES GCM usa la stessa chiave per cifrare e autenticare i dati, riducendo la complessità e il costo computazionale.
\end{itemize}
AES GCM è adatto per cifrare dati che richiedono accesso casuale, cioè che vengono letti o scritti in modo non sequenziale, o che vengono trasmessi su una rete.

\begin{figure}[h]
    \centering
    \includesvg[width=0.7\textwidth]{capitoli/risorse/AES_GCM_diagram.svg}
    \caption{Architettura AES GCM}
    \label{fig:AES-GCM-diagramma}
\end{figure}

\subsubsection{Cos'è il polinomio di Galois}
Il \textbf{polinomio di Galois} è un concetto fondamentale nella teoria dei campi e nell'algebra astratta. Inizialmente introdotto dal matematico Évariste Galois, il polinomio di Galois è associato a ogni estensione di campo finito.
Sia $F$ un campo e $E$ un'estensione di $F$. Se $E$ è un'estensione di Galois, allora il polinomio minimo di un elemento $a \in E$ su $F$ è detto \textbf{polinomio di Galois}.
Ad esempio, consideriamo l'estensione $E = \mathbb{Q}(\sqrt{2})$ su $F = \mathbb{Q}$, dove $\sqrt{2}$ è un numero irrazionale. Il polinomio minimo di $\sqrt{2}$ su $\mathbb{Q}$ è $x^2 - 2$. Questo è anche il polinomio di Galois per l'estensione $E/\mathbb{Q}$.
La teoria dei polinomi di Galois è profonda e ha numerose applicazioni in diverse aree della matematica, come la teoria dei numeri, la crittografia e la teoria dei codici.
Nel nostro caso il polinomio di Galois è utilizzato per generare un'autenticazione di messaggio durante il processo di cifratura e autenticazione dei dati. Questo polinomio è parte integrante dell'algoritmo di autenticazione del messaggio in GCM, che fornisce un modo efficiente per verificare l'integrità dei dati.
Il polinomio di Galois viene utilizzato per moltiplicare i blocchi di dati cifrati con una sequenza di valori di autenticazione derivati da una chiave di autenticazione. Questa operazione produce un'autenticazione dei dati, o MAC, che è inclusa insieme ai dati cifrati per consentire la verifica dell' integrità dei dati.

\subsubsection{AES CTR}
AES CTR (Counter Mode) è un modo di operazione che utilizza un contatore per cifrare i dati. In questo modo, ogni blocco di testo in chiaro viene sottoposto a un’operazione di XOR con il risultato della cifratura di un contatore, che è un numero che incrementa ad ogni blocco. Il contatore viene inizializzato con un nonce (numero usato una sola volta) e un IV. AES CTR incrementa un contatore per ogni blocco successivo cifrato. Ad esempio, se un’applicazione cifra una stringa venti volte di seguito, allora il risultato di ogni iterazione sarà diverso. Questo perché il contatore viene incrementato.

AES CTR ha diversi vantaggi:
\begin{itemize}
    \item \textbf{Parallelismo}: AES CTR è parallelizzabile, cioè è possibile cifrare o decifrare più blocchi contemporaneamente, perché ogni blocco dipende solo dal contatore. Questo aumenta le prestazioni, soprattutto con dati di grandi dimensioni.
    \item \textbf{Padding}: AES CTR non richiede padding. Normalmente, un algoritmo di cifratura a blocchi (AES, Blowfish, DES, RC2, ecc.) emette un output cifrato che è un multiplo della dimensione del blocco (16 byte per AES come esempio). Con la modalità CTR, il numero di byte di output è esattamente uguale al numero di byte di input, quindi non è richiesto alcun padding.
\end{itemize}
Questa modalità ha un grande svantaggio; a differenza di AES GCM, AES CTR non fornisce autenticazione. Questo significa che non è possibile verificare se il testo cifrato è stato alterato o manipolato da un attaccante. Per questa ragione, AES CTR è spesso utilizzato in combinazione con un codice di autenticazione del messaggio, come HMAC, per garantire sia la confidenzialità che l’integrità dei dati, tuttavia in TCFS non è stata implementata questa funzionalità. AES CTR è adatto per cifrare dati che richiedono accesso casuale, cioè che vengono letti o scritti in modo non sequenziale, o che vengono trasmessi su una rete.

\begin{figure}[H]
    \centering
    \includesvg[width=0.7\textwidth]{capitoli/risorse/CTR.svg}
    \caption{Architettura AES CTR}
    \label{fig:AES-CTR-diagramma}
\end{figure}

\subsubsection{Conclusione}
AES CBC, GCM e CTR sono tre modi di operazione di AES che presentano differenze significative in termini di sicurezza, prestazioni e applicazioni. AES GCM è generalmente più sicuro e più veloce di AES CBC, ma richiede una gestione accurata del nonce e dell’IV, che non devono essere riutilizzati per la stessa chiave. AES CBC è più semplice da implementare, ma non è parallelizzabile. AES CTR oltre a essere marginalmente più veloce di GCM offre prestazioni elevate e parallelizzazione, ma richiede una gestione accurata del contatore. Per questi motivi, sia AES GCM che AES CTR sono attualmente disponibili per la cifratura dei file, mentre AES CBC è utilizzato per i nomi delle directory. La scelta tra i tre dipenderà dalle specifiche esigenze dell’utente in termini di sicurezza, prestazioni e gestione delle chiavi.

\subsubsection{Sviluppi futuri e limitazioni attuali}
Attualmente TCFS non savla un IV per la cifratura dei nomi delle directory, questa limitazione è di caratterizzazione temporale, cioè nonostante lo sviluppo del progetto sia durato diversi mesi non si è riusciti a implementare questa funzionalità. Ci si propone quindi di implementarla quanto prima possibile.

\subsubsection{Valutazioni personali}
Attraverso l’analisi della situazione attuale si ricava che i benefici dovuti alla diversificazione degli algoritmi nei precedenti due scenari, si intende la cifratura dei nomi e la cifratura dei file stessi, venga surclassata dalla complessità manutentiva maggiorata.